%%
%% IntentCap extended abstract
%%

\documentclass[sigconf,nonacm]{acmart}

\setcopyright{none}
\settopmatter{printacmref=false,printccs=false}
\renewcommand\footnotetextcopyrightpermission[1]{}
\pagestyle{plain}

\AtBeginDocument{%
  \providecommand\BibTeX{{%
    Bib\TeX}}}

\emergencystretch=1.2em
\usepackage{microtype}
\usepackage{xspace}
\usepackage{tikz}
\usepackage{booktabs}
\usetikzlibrary{arrows.meta,positioning,fit,backgrounds}
\setlength{\bibsep}{0pt}

\setlength{\textfloatsep}{5pt plus 2pt minus 2pt}
\setlength{\dbltextfloatsep}{6pt plus 2pt minus 2pt}
\setlength{\floatsep}{6pt plus 2pt minus 2pt}
\setlength{\intextsep}{6pt plus 2pt minus 2pt}
\setlength{\abovecaptionskip}{3pt}

\newcommand{\sys}{\textsc{IntentCap}\xspace}

\begin{document}

\title{LLM Agent Capabilities Should Follow Task Intent and Context Source}
% LLM 智能体的能力应该跟随任务意图和上下文来源

\author{Yusheng Zheng$^1$, Wenhui Zhang$^2$, Yu Mao$^3$}
\affiliation{%
  \institution{$^1$UC Santa Cruz, USA \qquad $^2$Roblox, USA \qquad $^3$Bytedance, China}
  \city{}\country{}}
\email{yzhen165@ucsc.edu, wenhuizhang.psu@gmail.com, mynotwo@126.com}

\begin{abstract}
LLM agents take real actions---executing code, modifying files, calling services, delegating tasks---driven by context sources: user requests, tool results, documents, shell outputs, Skill and MCP instructions, memory.
% LLM 智能体会执行真实操作——运行代码、修改文件、调用服务、委托任务——由上下文来源 (context source) 例如用户请求、工具结果、文档、Shell 输出、Skill 与 MCP 指令和记忆所驱动。
Unlike traditional systems, where capability is predefined, the least-privilege capability an agent needs is dynamic, depending on its task intent: what it wants to do and how.
% 与能力是确定性且预定义的传统系统不同，智能体所需的最小权限能力是动态的，取决于智能体对任务的意图：它想做什么以及它计划如何完成。
This creates a security and safety challenge: all inputs enter one shared planning channel with equal influence, by adversarial injection or accidental scope widening, yet a user's explicit request and a document's extracted text carry different trust and should not share authority to choose a destination or widen access.
% 这带来了安全与可靠性挑战：所有输入以同等影响力进入同一个共享规划通道，无论是通过对抗性注入还是意外的范围扩大，但用户的显式请求和文档的提取文本有不同的信任等级，不应拥有同样的权限来选择目标或扩大访问范围。
Current defenses control allowed operations, not which inputs influence which decisions.
% 当前防御控制哪些操作被允许，但不控制哪些输入可以影响哪些决策。

We argue that agent capabilities should be scoped to the current task intent, not a sandbox or session lifetime, and that no single context source defines a complete capability.
% 我们认为智能体的能力应该限定在当前任务意图的范围内，而不是沙箱或会话生命周期，且没有任何单一上下文来源能定义完整的能力。
\sys composes capabilities from four such sources---user intent, workflow instructions, tool schemas, runtime environment---with field-level ownership and monotonic narrowing.
% IntentCap 从四个这样的来源——用户意图、工作流指令、工具模式和运行时环境——以字段级所有权和单调缩减来组合能力。
Each source contributes specific fields, none can fill another's, and the lease only narrows the user's authorized authority, never widens it.
% 每个来源贡献特定字段，任何来源都不能填充另一个来源的字段，由此产生的租约只能缩小用户授权的权限，不能扩大。
\sys uses an LLM to generate short-lived leases from these sources, validated by a deterministic checker before any side effect commits and enforced by tool- and OS-level information flow policies.
% IntentCap 使用 LLM 从这些来源生成短生命周期的能力租约，由确定性检查器在任何副作用提交之前验证，并通过工具级和 OS 级信息流策略执行。
Evaluation shows \sys blocks tested violations without rejecting benign actions, each source boundary is independently necessary, and the checker generalizes across tool, execution, placement, and delegation boundaries.
% 初步评估表明 IntentCap 阻断测试违规而不拒绝良性操作，每个来源边界都是独立必要的，且检查器可跨工具、执行、放置和委托边界推广。
\end{abstract}

\maketitle

\section{Background and Motivation}
% 背景与动机

LLM agents combine tools, documents, workflow instructions, and delegated subtasks to do real work.
% LLM 智能体结合工具、文档、工作流指令和委托的子任务来执行实际工作。
An agent can load Skills carrying workflow instructions, references, and scripts~\cite{openai-skills}, connect to MCP servers exposing resources, prompts, and tools~\cite{mcp-spec}, run local commands, and spawn subagents.
% 现代智能体可以加载包含工作流指令、引用和可选脚本的 Skill~\cite{openai-skills}，连接暴露资源、提示和工具的 MCP 服务器~\cite{mcp-spec}，运行本地命令，以及生成子智能体。
Each input is useful precisely because it affects future behavior---the same property that creates a security and safety problem.
% 同样的特性带来了安全与可靠性问题。
In an agent, any input can influence any action, through adversarial injection~\cite{greshake-injection} or accidental scope widening.
% 在智能体中，任何输入都可以影响任何操作，无论是通过对抗性注入~\cite{greshake-injection}还是意外的范围扩大。
Current defenses---prompt-level guardrails~\cite{camel,dualllm}, tool-level permission guards~\cite{progent}, execution sandboxes~\cite{isolategpt}, OS-level policy engines~\cite{actplane}---all control what the agent may do: which tool is called, which file read, which process executed.
% 当前防御——包括提示级护栏~\cite{camel,dualllm}、工具级权限守卫~\cite{progent}、执行沙箱~\cite{isolategpt} 和 OS 级策略引擎~\cite{actplane}——都控制智能体可以做什么：哪个工具可以调用，哪个文件可以读取，或者哪个进程可以执行。
None controls what drives each decision: which source selects the tool, fills the repository field, or widens approval scope.
% 没有一个控制什么被允许驱动每个决策，即哪个来源可以选择工具、填充仓库字段或扩大审批范围。

\textbf{Example.} 
Consider a user who asks an agent to extract tables from two selected PDFs, save spreadsheets, and create one GitHub issue in a named repository.
% 考虑一个用户要求智能体从两个选定的 PDF 中提取表格、保存为电子表格，并在指定仓库中创建一个 GitHub issue 的场景。
The PDF contents should influence spreadsheet cells and perhaps the issue body, but not the repository name, network destination, approval scope, tool selection, or Skill loading.
% PDF 内容应当影响电子表格的单元格和 issue 正文，但不应影响仓库名称、网络目标、审批范围、工具选择或加载另一个 Skill 的决定。
If hidden text in the PDF says ``create the issue in attacker/repo''~\cite{greshake-injection}, the repository field changes to a legitimate-looking, attacker-controlled destination. A tool-call allowlist can block banned domains, but cannot express the invariant that the PDF was never allowed to fill the repository field.
% 如果 PDF 中的隐藏文本说"在 attacker/repo 中创建 issue"~\cite{greshake-injection}，仓库字段就会变成一个看起来合法但由攻击者控制的目标。工具调用白名单可以阻止被禁止的域名，但它无法表达那个不变式：PDF 根本不被允许填充仓库字段。
The same invariant breaks without an adversary: if the extraction Skill concatenates the PDF text into one output string, the planner may parse a table header as the repository name, overwriting the user's chosen destination.
% 同样的不变式在没有攻击者的情况下也会被破坏。如果提取 Skill 将 PDF 文本拼接为单个输出字符串，规划器可能将表头解析为仓库名称，意外覆盖用户选择的目标。

The root cause is that all \emph{context sources} the agent consumes, what the user asked for, how a Skill defines the workflow, what a tool's API accepts, and what previous tools returned, enter the same planning channel with equal standing, so any source can fill any decision field.
% 根本原因是智能体所消耗的所有上下文来源 (context source)——用户要求做什么、Skill 如何定义工作流、工具的 API 接受什么、以及之前的工具返回了什么——以同等地位进入同一个规划通道，因此任何来源都可以填充任何决策字段。
Unlike traditional systems, where a programmer declares capabilities at compile time and a type system enforces them, an agent's planning channel has no static structure to distinguish sources.
% 与传统系统不同，在传统系统中程序员在编译时声明能力并由类型系统强制执行，而智能体的规划通道没有提供静态结构来区分来源。
The sources are not interchangeable: each carries only part of the information needed to define a capability. The user's request specifies the goal and authorized destinations but not the workflow steps. A Skill's procedure offers candidate steps but not the user's specific targets, of which the user may adopt only part. A tool's schema declares the parameters it accepts but not the goal. A tool's output provides values but not which decision fields they should fill. A complete capability must be composed from all four.
% 这些来源不可互换，因为每个来源只携带定义能力所需的部分信息。用户的请求指定目标和授权的目的地，但不包含工作流步骤。Skill 的流程提供候选步骤，但不包含用户的具体目标，且用户可能只采用其中一部分。工具的模式声明它接受什么参数，但不包含目标。工具的输出提供值，但不指定这些值应该填充哪些决策字段。完整的能力必须从四者组合而成。
Securing agents therefore requires an information flow policy controlling not only what operations are allowed, but which source may flow into which decision field.
% 因此，保护智能体的安全需要一种信息流策略 (information flow policy)，不仅控制哪些操作被允许，还控制哪个来源可以流入哪个决策字段。

\section{Design}
\label{sec:design}
% 设计

Our key insight is that agent capabilities must be composed at runtime from multiple context sources with field-level ownership, because no single source carries enough authority to define a complete capability, and no programmer exists to declare the composition statically.
% 我们的关键洞察是智能体的能力必须在运行时从多个上下文来源以字段级所有权进行组合，因为没有单一来源能定义完整的能力，也没有程序员来静态声明这种组合。
\sys realizes this by scoping each capability to the current task intent and composing it from four sources, each contributing specific fields with priority, so none can fill another's.
% IntentCap 通过将每个能力限定在当前任务意图的范围内，并从四个来源组合，每个来源按优先级贡献特定字段，使得任何来源都不能填充另一个来源的字段来实现这一点。

The design targets four properties: (1) field-level ownership, where each decision field is owned by exactly one context source; (2) monotonic narrowing, where leases can only narrow the user's authorized authority, never widen it; (3) the LLM stays outside the trusted computing base, with a deterministic checker making all accept/deny decisions; and (4) auditable leases, where every accept or deny is recorded with its provenance chain.
% 设计目标四个属性：(1) 字段级所有权——每个决策字段由恰好一个上下文来源拥有；(2) 单调缩减——租约只能缩小用户授权的权限，不能扩大；(3) LLM 位于可信计算基之外——确定性检查器做出所有接受/拒绝决策；(4) 可审计租约——每个接受或拒绝都记录来源链。

\textbf{Threat model and TCB.}
The adversary controls runtime-environment content, documents, tool outputs, web data, delegated-agent messages, and may plant instructions in Skill or MCP text; it does not control the user's structured selections or the enforcement layer.
% 威胁模型与 TCB。攻击者控制运行时环境的内容——文档、工具输出、网络数据、委托智能体的消息——并可能在 Skill 或 MCP 文本中植入指令；它不能控制用户的结构化选择或执行层。
The LLM planner and compiler are untrusted and arbitrarily manipulable: they propose but never decide.
% LLM 规划器和 LLM 辅助编译器不受信任且可能被任意操纵：它们只提出建议，从不做决定。
The TCB (dashed box in Figure~\ref{fig:arch}) is the intent issuer, source labeler, deterministic checker, and enforcement layer (hooks, MCP gateway, OS enforcer).
% TCB——即图~\ref{fig:arch} 中的虚线框——是意图签发器、来源标注器、确定性检查器和执行层（钩子、MCP 网关、OS 执行器）。

\begin{figure}[t]
\centering
\resizebox{\columnwidth}{!}{%
\begin{tikzpicture}[
  x=1cm, y=1cm,
  box/.style={draw, rounded corners=2pt, align=center, minimum height=0.8cm,
              text width=2.5cm, fill=white, font=\large, inner sep=2pt},
  trusted/.style={box, fill=green!8},
  untrusted/.style={box, fill=red!7},
  runtime/.style={box, fill=blue!7, text width=2.3cm},
  lbl/.style={font=\large, inner sep=1.5pt},
  arr/.style={-Latex, semithick}
]
% Trusted chain: one row, so the dashed TCB box stays contiguous.
\node[trusted] (issuer)  at (0,0)   {Intent issuer};
\node[trusted] (labeler) at (3.8,0) {Source labeler};
\node[trusted] (checker) at (7.6,0) {Deterministic checker};
% Untrusted compiler: above the checker, outside the TCB.
\node[untrusted] (compiler) at (7.6,1.55) {LLM-assisted compiler};
% Enforcement layer.
\node[runtime] (adapters) at (7.6,-1.55)  {Enforcement hooks};
\node[runtime] (gateways) at (4.6,-1.55)  {MCP gateway};
\node[runtime] (deleg)    at (10.6,-1.55) {OS-level enforcer};

\draw[arr] (issuer) -- (labeler);
\draw[arr] (labeler.north) |- (compiler.west);
\draw[arr] ([xshift=6mm]compiler.south) --
      node[lbl, right] {candidates} ([xshift=6mm]checker.north);
\draw[arr, dashed] ([xshift=-6mm]checker.north) --
      node[lbl, left] {denial} ([xshift=-6mm]compiler.south);
\draw[arr] (checker) -- node[lbl, right] {leases} (adapters);
\draw[arr] (adapters) -- (gateways);
\draw[arr] (adapters) -- (deleg);

\begin{scope}[on background layer]
\node[draw, dashed, rounded corners=3pt, fit=(issuer)(labeler)(checker),
      inner sep=0.16cm, label={[lbl]below:TCB}] {};
\end{scope}
\end{tikzpicture}%
}
\caption{\sys architecture: the untrusted compiler proposes leases, the TCB checker decides.}
% \Description{The intent issuer, source labeler, and deterministic checker form the trusted computing base; an untrusted LLM-assisted compiler proposes candidate leases, and accepted leases are lowered to the MCP gateway, enforcement hooks, and OS-level enforcer.}
\label{fig:arch}
\end{figure}

Figure~\ref{fig:arch} shows the pipeline.
% 图~\ref{fig:arch} 展示了该流水线。
A source labeler classifies each input by position.
% 来源标注器按位置分类每个输入。
User messages are always user intent, tool responses always runtime environment, and Skill/MCP metadata is labeled at load time.
% 用户消息始终是用户意图，工具响应始终是运行时环境，Skill/MCP 元数据在加载时标注。
Inputs with ambiguous provenance (e.g., user-pasted documents) require additional classification, which the prototype handles conservatively by denying them authority fields.
% 来源模糊的输入（如用户粘贴的文档）需要额外分类，当前原型保守地拒绝模糊来源的权限字段。
At the OS level, this labeling maps to process isolation and IPC channel tagging, rooting each source's provenance in the process that produced it.
% 在 OS 层面，这种标注映射到进程级隔离和 IPC 通道标记，每个来源的出处都根植于产生它的进程。
The intent issuer extracts the user's maximum authority from structured fields in the user message: selected files, named destinations, and explicit approvals.
% 意图签发器从用户消息的结构化字段中提取用户的最大权限，例如选定文件、命名目标和显式批准。
The compiler takes this structured intent and the active schemas and emits candidate leases whose fields are populated from the labeled sources.
% 编译器以该结构化意图和活跃模式为输入，输出候选租约，其字段从已标注的来源填充。
Candidate leases can only narrow the user's authority, never widen it.
% 候选租约只能缩小用户的权限，不能扩大。
The compiler may select a subset of authorized destinations or request a shorter budget, but cannot add destinations, raise approval scope, or grant fields the user intent does not cover.
% 编译器可以选择授权目标的子集或请求更短的预算，但不能添加目标、提高审批范围或授予用户意图未涵盖的字段。


A deterministic checker enforces this monotonic narrowing before any side effect, context placement, or delegation handoff commits.
% 确定性检查器在任何副作用、上下文放置或委托交接提交之前强制执行这种单调缩减。
Each delegation further attenuates, so a subtask receives only capabilities narrower than its parent.
% 每次委托进一步衰减，子任务只接收比其父任务更窄的能力。
The checker exposes a single transition, $\mathsf{check\_and\_consume}(e,\mathit{lease},\mathit{proofs},\sigma) \rightarrow \mathsf{allow}(\sigma') \mid \mathsf{deny}$, that every adapter must call before a side effect commits.
% 检查器暴露单一转换——$\mathsf{check\_and\_consume}(e,\mathit{lease},\mathit{proofs},\sigma) \rightarrow \mathsf{allow}(\sigma') \mid \mathsf{deny}$——适配器必须在任何副作用提交之前调用。
Enforcement has two layers. At the tool level, boundary adapters intercept MCP calls and validate that each argument field comes from its owner source before execution, as Fides~\cite{fides} applies taint tracking at the planner level. At the OS level, accepted leases compile to ActPlane~\cite{actplane} file, exec, and network policies enforced via eBPF.
% 执行在两个层次运作。在工具层，边界适配器拦截 MCP 调用并在调用执行前验证每个参数字段来自其所有者来源，类似于 Fides~\cite{fides} 在规划器层面应用污点追踪。在 OS 层，被接受的租约编译为通过 eBPF 执行的 ActPlane~\cite{actplane} 文件、执行和网络策略。

A field is a named slot in a capability lease (e.g., destination, approval scope, tool name, argument value) whose value must come from exactly one source.
% 字段是能力租约中的一个命名槽位（如目标、审批范围、工具名、参数值），其值必须来自恰好一个来源。
The four context sources each contribute specific fields.
% 四个上下文来源各自贡献特定字段。
\emph{User intent} supplies the goal, selected objects, authorized destinations, and approvals.
% \emph{用户意图}提供目标、选定对象、授权目标和批准。
\emph{Workflow instructions} (system policy, Skill procedures, or manuals) supply procedural scope.
% \emph{工作流指令}（系统策略、Skill 流程或手册）提供过程范围。
\emph{Tool schemas} (MCP definitions, registry entries, or command descriptors) supply the callable interface and credential scope.
% \emph{工具模式}（MCP 定义、注册表条目或命令描述符）提供可调用接口和凭证范围。
\emph{Runtime environment} (tool results, file state, script outputs) supplies observed values.
% \emph{运行时环境}（工具结果、文件状态、脚本输出）提供观测值。
No source can fill another's fields.
% 任何来源都不能填充另一个来源的字段。
A tool schema cannot authorize a destination, runtime environment cannot supply procedural authority, and a Skill cannot widen approval scope.
% 工具模式不能授权目标，运行时环境不能提供过程权限，Skill 不能扩大审批范围。
Together these ownership rules form an information flow policy for agent decisions: which information may flow from which source into which field of a capability lease.
% 这些所有权规则共同构成了智能体决策的信息流策略：它们指定哪些信息可以从哪个来源流入能力租约的哪个字段。

A capability lease binds these field contributions to a specific operation, object, arguments, budget, expiry, and delegation depth.
% 能力租约将这些字段贡献绑定到特定的操作、对象、参数、预算、过期时间和委托深度。
Leases are scoped to the current task intent, consumed on use, and cannot widen themselves.
% 租约限定在当前任务意图的范围内，使用即消耗，且不能自行扩大。
An issue lease expires after first use, and a delegated subtask receives only capabilities narrower than its parent's.
% 一次性 issue 租约在首次使用后即过期，被委托的子任务只接收比其父任务更窄的能力。
The checker validates all fields atomically in a single transition before the side effect commits.
% 检查器在副作用提交之前，在单次转换中原子地验证所有字段。

In the PDF task above, user intent supplies the two selected files and repository \texttt{org/proj}, the Skill the extract--save--report procedure, the issue tool's schema the callable interface, and runtime results the table values.
% 工作示例。在上述 PDF 任务中，用户意图提供两个选定文件和仓库 \texttt{org/proj}，Skill 提供提取—保存—报告的流程，issue 工具的模式提供可调用接口，运行时结果提供表格值。
The compiler proposes a lease for \texttt{create\_issue}: the repository field copies \texttt{org/proj} with a witness to the user-intent span, the issue body is a data field synthesized from the PDFs, budget one issue, expiring on first use.
% 编译器为 \texttt{create\_issue} 提出一个租约：仓库字段复制 \texttt{org/proj} 并附带指向用户意图片段的见证，issue 正文是从 PDF 合成的数据字段，预算为一个 issue，首次使用即过期。
The checker verifies witness and narrowing, and commits.
% 检查器验证见证与缩减后提交。
An injected candidate naming \texttt{attacker/repo} has no witnessing user-intent span, so the checker denies it before the call executes.
% 注入产生的指定 \texttt{attacker/repo} 的候选租约没有任何用户意图片段为其作证，检查器在调用执行前拒绝该租约。

\section{Preliminary Evaluation}
% 初步评估

\textbf{RQ1: Does \sys block violations without rejecting benign actions?}
% \textbf{RQ1：IntentCap 能否阻断违规而不拒绝良性操作？}
\sys produces 0 unsafe accepts across 3,746 security-sensitive events replayed from AgentDojo, MCPTox, InjecAgent, and tau2-bench~\cite{agentdojo,mcptox,injecagent,tau2bench}.
% IntentCap 在从 AgentDojo、MCPTox、InjecAgent 和 tau2-bench 重放的 3,746 个安全敏感 (security-sensitive) 事件中产生 0 个不安全接受~\cite{agentdojo,mcptox,injecagent,tau2bench}。
For utility, \sys covers all 3,813 benign reference actions in the tau2-style proxy and passes 2,554 of 2,556 applicable ground-truth checks.
% 在可用性方面，IntentCap 覆盖了 tau2 风格代理中全部 3,813 个良性参考操作，并通过了 2,556 个适用的真实标签 (ground-truth) 检查中的 2,554 个。

\textbf{RQ2: Is the four-source partition a safety requirement or a naming choice?}
% \textbf{RQ2：四来源划分是安全要求还是命名选择？}
Collapsing all four sources into one generic trusted context causes 3,593 false accepts among the 3,823 events the checker should deny (94\%).
% 将四个来源合并为一个通用受信任上下文导致 3,823 个检查器应拒绝事件中出现 3,593 个错误接受（94\%）。
Pairwise collapses open distinct holes: tool$\rightarrow$agent falsely accepts 1,928 events (50\%), env$\rightarrow$agent 1,663 (43\%), env$\rightarrow$tool 1,662 (43\%).
% 成对合并各自打开不同的漏洞：tool$\rightarrow$agent 错误接受 1,928 个事件（50\%），env$\rightarrow$agent 1,663（43\%），env$\rightarrow$tool 1,662（43\%）。
On 7 crafted multi-step workflows, a policy DSL checking predicates without field ownership falsely accepts 7/7, and splitting state across independent guards 5/7.
% 在 7 个手工构造的多步工作流场景上，一个不检查字段所有权只检查谓词的策略 DSL 错误接受了 7/7，将状态拆分到各独立守卫则错误接受了 5/7。

\textbf{RQ3: Does the same commit API work beyond tool calls?}
% \textbf{RQ3：同一提交 API 能否超越工具调用生效？}
\texttt{check\_and\_consume} generalizes across five enforcement points: local env side effects, context placement, Skill instruction slots, delegation handoffs, and an OS monitor replay backend.
% \texttt{check\_and\_consume} 转换可推广到五种执行点，包括本地环境副作用、上下文放置、Skill 指令槽、委托交接和 OS 监视器风格的重放后端。
Across 38 checks, \sys allows 17 authorized effects, blocks 21 violations, with 0 unsafe executions or placements and 0 checker/monitor mismatches.
% 在 38 个执行检查中，IntentCap 允许了 17 个授权效果并阻断了 21 个违规，0 个不安全执行或放置，0 个检查器/监视器不匹配。

\textbf{RQ4: Are leases auditable authority objects?}
% \textbf{RQ4：租约是否为可审计的权限对象？}
Across 24 hand-labeled leases from InjecAgent, MCPTox, and tau2, all \sys policies match the labeled scope; every non-\sys baseline over-grants on at least some of the 144 scored policy entries.
% 在覆盖 InjecAgent、MCPTox 和 tau2 样本的 24 个人工标注的租约上，所有 IntentCap 策略与标注的权限范围完全匹配；每个非 IntentCap 基线在 144 个评分策略条目中的至少部分条目上过度授权。

\textbf{Limitations and future work.}
Unlike Capsicum~\cite{capsicum} file descriptors or seccomp filters, an agent's capability cannot be declared at compile time~\cite{dennis-vanhorn,saltzer-schroeder,denning-lattice,confused-deputy}; against defenses that separate control from data~\cite{camel}, restrict operations~\cite{progent,isolategpt}, or constrain flow, tools, the OS, and packages~\cite{fides,eim-bpftime,actplane,skillguard}, \sys binds each decision field to its single authorized source.
% 与 Capsicum 文件描述符或 seccomp 过滤器不同，智能体的能力无法在编译时声明；相对于将控制与数据分离、限制可执行操作、或约束信息流/工具/OS/包的防御，IntentCap 将每个决策字段绑定到其唯一授权来源。
Copy-witness provenance is conservative: cross-source synthesis is confined to data fields or user-authorized selectors.
% 复制见证出处是保守的：跨来源合成被限制在数据字段或用户授权的选择器内。
Meet-based composition, end-to-end attack success, enforcement latency, and adaptive attacks on labeling and lease generation remain future work.
% 基于 meet 的组合、端到端攻击成功率、执行延迟，以及针对标注和租约生成的自适应攻击仍是未来工作。

\bibliographystyle{ACM-Reference-Format}
\bibliography{references}

\end{document}
